\chapter{R\'esultats obtenus}
\label{chap:resultats}

Ce chapitre pr\'esente les r\'esultats techniques obtenus lors du d\'eveloppement de
notre application~: l'architecture mise en place, les mesures de s\'ecurit\'e
appliqu\'ees, ainsi que certaines interfaces illustrant le fonctionnement de la
plateforme.

\section{R\'esultats}

\subsection{Pr\'esentation de quelques interfaces de l'application}

% TODO : remplacer chaque \TODOFigure par une vraie capture d'\'ecran une fois
% le frontend d\'evelopp\'e (page d'accueil, tableau de bord, nouvelle analyse,
% r\'esultat d'analyse, historique, administration...).

\subsubsection{Page d'accueil}
\textit{[\`A compl\'eter]}
\TODOFigure{page-accueil.png}{Page d'accueil}{fig:accueil}

\subsubsection{Tableau de bord}
\textit{[\`A compl\'eter]}
\TODOFigure{tableau-de-bord.png}{Tableau de bord}{fig:dashboard}

\subsubsection{R\'esultat d'une analyse}
\textit{[\`A compl\'eter]}
\TODOFigure{resultat-analyse.png}{R\'esultat d'une analyse}{fig:resultat}

\subsection{Syst\`eme de gestion de base de donn\'ees}

Pour ce projet, nous avons utilis\'e \textbf{MySQL} (MariaDB) comme syst\`eme de
gestion de base de donn\'ees relationnelle, int\'egr\'e via XAMPP. La base de donn\'ees
compte quatre tables principales~: \texttt{utilisateurs}, \texttt{analyses},
\texttt{controles} et \texttt{parametres\_score}, refl\'etant directement le
diagramme de classes pr\'esent\'e au chapitre pr\'ec\'edent. L'utilisation de XAMPP
facilite le d\'eveloppement en local en fournissant un environnement complet
incluant Apache, PHP et MySQL. Laravel interagit avec MySQL gr\^ace \`a
\textbf{Eloquent}, son ORM int\'egr\'e, qui permet de manipuler les donn\'ees de
mani\`ere simple, s\'ecuris\'ee et orient\'ee objet.

\subsection{Protocole de s\'ecurit\'e}

La s\'ecurit\'e a \'et\'e au c\oe ur de notre d\'emarche de d\'eveloppement, ce qui est
particuli\`erement coh\'erent avec la nature m\^eme du projet.

\paragraph{Authentification} L'authentification repose sur \textbf{Laravel
Sanctum} en mode SPA (authentification par cookie de session), adapt\'ee \`a la
s\'eparation frontend React / backend Laravel. Les mots de passe sont stock\'es
hach\'es (\texttt{bcrypt} via \texttt{Hash::make}), jamais en clair.

\paragraph{Autorisation} Chaque route sensible de l'API v\'erifie que
l'utilisateur authentifi\'e est bien propri\'etaire de la ressource demand\'ee (par
exemple, un utilisateur ne peut consulter le d\'etail d'une analyse que si elle
lui appartient, sauf s'il est administrateur), avec un rejet explicite (code
HTTP 403) sinon.

\paragraph{Validation des entr\'ees} Toutes les donn\'ees entrantes (URL \`a
analyser, formulaires d'inscription/connexion) sont valid\'ees c\^ot\'e serveur via
le validateur int\'egr\'e de Laravel avant tout traitement.

\subsection{Protection des donn\'ees et communication s\'ecuris\'ee}

\begin{itemize}
  \item Les injections SQL sont \'evit\'ees gr\^ace \`a l'ORM Eloquent, qui prot\`ege
    automatiquement les requ\^etes contre ces attaques.
  \item Le moteur d'analyse lui-m\^eme reste strictement \textbf{non intrusif}~: il
    se limite \`a observer les r\'eponses HTTP publiques d'un site (en-t\^etes,
    certificat) sans jamais tenter d'exploiter une faille.
  \item % TODO : une fois d\'eploy\'ee, pr\'eciser ici que les \'echanges entre le
        % frontend et l'API se font via HTTPS en production.
        \textit{[\`A compl\'eter apr\`es d\'eploiement~: confirmation du HTTPS en production]}
\end{itemize}

\section{Discussion}

\subsection{Par rapport aux objectifs fix\'es}

Les objectifs sp\'ecifiques fix\'es en introduction sont, \`a ce stade,
globalement atteints~: la plateforme authentifie ses utilisateurs (par mot de
passe ou via Google), analyse un site sur les sept crit\`eres retenus, calcule
un score et un niveau de risque, g\'en\`ere des recommandations, et conserve un
historique consultable. La principale valeur ajout\'ee par rapport aux outils
existants cit\'es en introduction (Mozilla Observatory, SSL Labs) tient moins
aux contr\^oles eux-m\^emes -- volontairement plus restreints -- qu'\`a leur
restitution~: un score unique plut\^ot que plusieurs rapports \'epars, et des
descriptions r\'edig\'ees en fran\c{c}ais courant plut\^ot qu'en jargon technique
anglophone. Ce dernier point n'\'etait pas acquis d\`es le d\'epart~: les
premi\`eres versions du moteur produisaient des constats bruts (noms
d'en-t\^etes HTTP, versions de protocole) qui se sont r\'ev\'el\'es tout aussi
incompr\'ehensibles pour un non-technicien que ceux des outils existants. La
reformulation compl\`ete des descriptions et recommandations a \'et\'e n\'ecessaire
pour r\'epondre effectivement \`a l'objectif de p\'edagogie.

\subsection{Limites}

Plusieurs limites doivent \^etre soulign\'ees. D'abord, le score refl\`ete un
\'etat \`a un instant T~: une analyse est une photographie, pas une surveillance
continue -- un site conforme au moment de l'analyse peut se d\'egrader ensuite
sans que l'utilisateur en soit inform\'e, hors nouvelle analyse manuelle.
Ensuite, les sept contr\^oles retenus, bien que repr\'esentatifs, ne constituent
pas un audit de s\'ecurit\'e exhaustif~: ils ne remplacent pas un test
d'intrusion men\'e par un professionnel, et se limitent volontairement \`a des
v\'erifications passives. La pond\'eration des crit\`eres dans le calcul du score,
bien que configurable par un administrateur, reste \'egalement une d\'ecision en
partie subjective~: deux baremes diff\'erents peuvent produire deux niveaux de
risque diff\'erents pour un m\^eme site.

Le traitement des analyses d\'ependant d'appels r\'eseau vers des serveurs
externes r\'eels, le r\'esultat peut en outre varier d'une ex\'ecution \`a
l'autre~: au cours des tests, une m\^eme URL a pu \^etre marqu\'ee \og~injoignable~\fg{}
lors d'un essai puis analys\'ee avec succ\`es quelques minutes plus tard,
simplement parce que le site cibl\'e a mis plus de huit secondes \`a r\'epondre la
premi\`ere fois. Ce comportement n'est pas un dysfonctionnement de la
plateforme, mais une contrainte inh\'erente \`a toute analyse d\'ependante du
r\'eseau, qu'il convient de garder \`a l'esprit dans l'interpr\'etation d'un
\'echec ponctuel.

Enfin, le traitement asynchrone des analyses (section~\ref{chap:programmation})
repose sur un worker de file d'attente qui doit tourner en continu~: en
environnement de d\'eveloppement local, l'oubli de d\'emarrer ce processus a \`a
plusieurs reprises laiss\'e des analyses bloqu\'ees ind\'efiniment au statut
\og~en cours~\fg, sans message d'erreur explicite pour l'utilisateur. Ce point
constitue une vigilance op\'erationnelle \`a documenter pour un futur
d\'eploiement en production, au-del\`a du seul code applicatif.

Ce chapitre a permis de pr\'esenter les principaux r\'esultats techniques du
projet ainsi qu'une lecture critique de leur port\'ee et de leurs limites.
